\Needspace{13\baselineskip}
\chapterrule
\section{把语言练习变成一道可计算的预测题}
\label{intro:20b-task-contract}

下一 Token 预测把已有文本变成训练答案，把新 Token 接回输入又能继续生成。
现在把这个任务写成一份\textbf{任务契约（Task Contract）}：
先确定输入、目标、输出形状和计分方式。\phantomsection\label{concept:task-contract}
模型结构随后遵守这份约定。

\subsection{一句话怎样拆成三道预测题？}

沿用五项教学词表：$1=\text{其他}$、$2=\text{春天}$、
$3=\text{。}$、$4=\text{花}$、$5=\text{开了}$。
本次句子使用其中四个 Token“春天／花／开了／。”，编号为
\[
(t_1,t_2,t_3,t_4)=(2,4,5,3).
\]
在实际模型中，Token 的单位包括单字、词的一部分和标点。

先看文字怎样变成题目。沿着原句从左向右读，\textbf{每读到一个 Token，
就停下来问：“紧接着的下一个 Token 是什么？”}原文中紧跟在后面的 Token，
就是这道题的正确答案，也叫\textbf{目标}。这句话可以提供三道题：

\begin{center}
\small
\tagpdfsetup{table/header-rows={1}}
\begin{tabularx}{\linewidth}{@{}p{22mm}Xp{35mm}@{}}
\toprule
预测位置 & 读到这里，可以利用的文字 & 正确答案（目标）\\
\midrule
1 & 春天 & 花（ID $4$）\\
2 & 春天、花 & 开了（ID $5$）\\
3 & 春天、花、开了 & 。\quad（ID $3$）\\
\bottomrule
\end{tabularx}
\end{center}

“位置”按输入 Token 的先后顺序编号。位置~2 就是读到第二个 Token“花”的时候；
此时已经读过“春天、花”，要预测的是它们后面的“开了”。
位置~3 读到“开了”，下一项是句号。句号在这次练习中用作最后一道题的答案。

按预测位置依次记录\textbf{当前读到的 Token}，得到输入“春天、花、开了”；
再按相同顺序记录\textbf{各位置要预测的下一 Token}，得到目标“花、开了、。”。用 ID 表示：
\[
\underbrace{[2,\ 4,\ 5]}_{\text{三个位置的输入 ID}}
\qquad
\underbrace{[4,\ 5,\ 3]}_{\text{逐位置的目标 ID}}.
\]
输入取原文的前三项，目标取各输入右边紧邻的一项，\textbf{按位置一一配对}。
例如，位置~2 的当前输入是“花”（ID $4$），目标是“开了”（ID $5$），
做预测时还可读取前面的“春天”；三个位置可利用的前文分别见表中的范围。

\subsection{模型接收什么，每个位置输出什么？}

把预测模型暂时看成一个待展开的盒子。输入是 $[2,4,5]$，
它要对三个位置各回答一次“下一 Token 是谁”。
每道题都有五个候选，所以每个位置输出五项分数，合起来是一个 $3\times5$ 的矩阵。
\textbf{每一行放一道题的五项分数}：第~1 行对应位置~1，评价“春天”后面接谁；
第~2 行对应位置~2，评价“春天、花”后面接谁；第~3 行对应位置~3，评价“春天、花、开了”后面接谁。
矩阵的三行与前表中的三个预测位置一一对应，五列依次对应五个词表 ID。

分数怎样变成概率、怎样与正确答案比较，将在下一章计算；
整数 ID 怎样变成可训练的向量，再由 Embedding 一章展开。
训练时，用每行目标 ID 指出的答案评价该行预测，再汇总成损失 $L$。
例如第一行评价“花”，第二行评价“开了”，第三行评价“。”。
目标分支进入评分步骤；模型做预测时，只能使用本题允许读取的输入。

\par\noindent\begin{minipage}{\linewidth}
\tagstructbegin{tag=Figure,alttext={下一 Token 预测的训练任务：输入 Token ID 二、四、五经过预测模型，得到三行五列的候选分数。预测与另一分支上的目标 ID 四、五、三共同进入评分步骤，得到损失 L。目标只进入评分步骤。}}
\begin{center}
\begin{tikzpicture}[>=Latex,node distance=6mm,
 n/.style={draw=linegray,rounded corners=2pt,fill=white,align=center,
 font=\sffamily\footnotesize,minimum height=12mm,minimum width=28mm,inner sep=1mm}]
\node[n,fill=paper] (input) {输入 Token ID\\ $[2,4,5]$};
\node[n,right=of input,minimum width=32mm] (model) {预测模型\\ 内部随后展开};
\node[n,right=of model] (probability) {五项候选分数\\ 共 $3\times5$};
\node[n,right=of probability,fill=linegray!15] (loss) {比较预测与目标\\ 得到损失 $L$};
\node[n,below=6mm of probability,fill=paper] (target) {目标 Token ID\\ $[4,5,3]$};
\foreach \a/\b in {input/model,model/probability,probability/loss}{\draw[->,flowcolor,thick] (\a.east)--(\b.west);}
\draw[->,flowcolor,thick] (target.east)-|(loss.south);
\end{tikzpicture}
\captionof{figure}{先确定预测题：模型输出三行候选分数，错开一位的目标用于评价预测}
\label{fig:token-ffn-output-pipeline}
\end{center}
\tagstructend
\end{minipage}\par

真实 GPT-2 small 的词表有 $50{,}257$ 项，仍沿用下一 Token 目标；
规模与配置见第~\ref*{intro:23a-from-toy-to-gpt}~章。

\subsection{每个位置能看到哪些输入？}

预测下一 Token 时，每个位置可以利用\textbf{当前位置及之前}的输入。
这就是前表中三种可见范围的来由：读到哪里，就用截至那里的文字预测下一项。
把三道题一起计算时，虽然输入 $[2,4,5]$ 已经包含“花”和“开了”，
位置~1 的预测仍只能使用“春天”，位置~2 的预测可以使用“春天、花”。
例如，位置~1 要预测“花”，若计算预测时已经读到了位置~2 的“花”，
就提前看到了答案，造成\textbf{信息泄漏}。
目标应当参与损失计算；预测模型怎样利用允许读取的前文，将从下一部分逐步展开。
下一章先解决评分问题：五项实数分数怎样变成概率，怎样衡量它们与正确答案的差距？

\chapterchecktasks{20b-task-contract}{理解与复算}{
\item \textbf{对齐：}完整序列 $[2,4,5,3]$，采用三个输入位置 $[2,4,5]$。列出逐位置的三个下一 Token 目标。\chapterinlinehint{在完整文本中找到每个输入 Token，取它右边紧邻的一项作目标。}
\item \textbf{形状：}词表大小为 $5$，这条示例的三个位置各输出五类分数。把这些分数组织成矩阵后是什么形状？\chapterinlinehint{一条位置对应一行，类别数对应列数。}
\item \textbf{找错：}同学让第一个位置直接读取第二个位置的“花”，再预测第一个位置的下一 Token。哪里发生信息泄漏？\chapterinlinehint{先写出第一个位置的目标，再检查它是否已出现在预测所用的输入中。}
}
